\chapter{Dust Mites}
\index{Dust Mites}

\section*{Definition and Overview}
Dust mites are microscopic creatures that live in house dust and thrive in warm, humid environments such as bedding, mattresses, carpets, and upholstered furniture. They feed on the flakes of skin shed by people and pets. Dust mites themselves do not bite or cause harm directly, but their droppings and body fragments contain proteins that trigger allergic reactions in many people, including asthma, allergic rhinitis, and eczema. This chapter explains how dust mites affect health and the practical steps that reduce exposure.

\section*{Epidemiology}
Dust mites are found in homes throughout the world and are one of the most common triggers of allergy. A large proportion of people with allergic rhinitis, asthma, or eczema are sensitised to dust mite allergens, and in many countries, where the climate is warm and humid in many regions, dust mite allergy is particularly common. Exposure is greatest in the bedroom, where people spend a third of their lives in close contact with mattresses and pillows.

\section*{Aetiology and Pathophysiology}
Dust mite allergy is a hypersensitivity reaction to proteins in mite faeces and body parts.
\begin{itemize}
    \item \textbf{The mites:} Dermatophagoides species are the most important; they flourish in warm, humid conditions and feed on shed skin flakes
    \item \textbf{The allergens:} proteins in mite droppings and bodies are inhaled when dust is disturbed
    \item \textbf{The allergic response:} in sensitised people, the immune system produces antibodies that trigger the release of histamine and other chemicals, causing inflammation
    \item \textbf{Effects:} inflammation of the nasal passages, airways, and skin produces the symptoms of hay fever, asthma, and eczema
    \item \textbf{Environmental factors:} warmth, humidity, and a plentiful food supply of skin flakes favour mite populations
\end{itemize}

\section*{Clinical Features}
Symptoms depend on where the allergic reaction occurs and are typically worst at night and in the morning.
\begin{itemize}
    \item Sneezing, runny or blocked nose, and itchy eyes, nose, and throat
    \item Wheezing, coughing, chest tightness, and breathlessness in people with asthma
    \item Itchy, dry, and inflamed skin in people with eczema
    \item Symptoms worse in the bedroom, on waking, or after making the bed
    \item Worsening of symptoms in humid seasons
    \item Fatigue from disrupted sleep caused by nasal blockage or asthma at night
\end{itemize}

\section*{Differential Diagnosis}
Allergy to dust mites may be confused with other causes of similar symptoms.
\begin{itemize}
    \item Seasonal hay fever from pollens, and allergies to other indoor allergens such as pet dander and mould
    \item Viral colds and other respiratory infections
    \item Non-allergic rhinitis, triggered by temperature, strong smells, or smoke
    \item Sinusitis and nasal polyps
    \item Eczema caused by irritants or other allergens
\end{itemize}

\section*{When to Seek Medical Care}
See a GP if allergy symptoms are persistent, troublesome, or interfering with sleep, school, or work, or if asthma is not well controlled. A doctor can arrange allergy testing and prescribe appropriate treatment.

\textbf{Contact your local emergency services for:}
\begin{itemize}
    \item Severe difficulty breathing or wheezing that does not improve with reliever medication
    \item Chest tightness that is severe or worsening rapidly
    \item Signs of a severe allergic reaction, including throat swelling, difficulty swallowing, or collapse
\end{itemize}

\section*{Diagnosis and Investigations}
Diagnosis combines the clinical history with allergy testing.
\begin{itemize}
    \item \textbf{History:} symptom pattern, time of day, season, and triggers
    \item \textbf{Allergy testing:} skin prick tests or blood tests for specific antibodies confirm sensitisation to dust mites
    \item \textbf{Peak flow testing:} for people with asthma, to assess airway function
    \item \textbf{Assessment of the home:} identifying the main sources of mite allergen in the bedroom
\end{itemize}

\section*{Management}
Treatment aims to control symptoms and reduce exposure.
\begin{itemize}
    \item \textbf{Allergen avoidance:} the most important measure, focusing on the bedroom
    \item \textbf{Medication:} antihistamines and nasal corticosteroid sprays for rhinitis; asthma preventers and relievers as prescribed
    \item \textbf{Immunotherapy:} allergy shots or tablets that gradually desensitise the immune system, effective for many people with dust mite allergy
    \item \textbf{Treatment of eczema:} moisturisers, topical corticosteroids, and other prescribed treatments
    \item \textbf{Cleaning strategy:} frequent vacuuming with a HEPA-filter vacuum, washing bedding weekly in hot water, and reducing soft furnishings
\end{itemize}

\section*{Complications}
\begin{itemize}
    \item Poorly controlled asthma, with increased risk of asthma attacks and hospital admission
    \item Chronic nasal congestion, which may lead to sinusitis and sleep disturbance
    \item Disrupted sleep and daytime fatigue
    \item Impaired concentration and performance at school or work
    \item Persistently troublesome eczema with risk of skin infection
\end{itemize}

\section*{Prognosis and Outcomes}
Dust mite allergy is a chronic condition, but symptoms are well controlled in most people with a combination of allergen reduction and medication. Immunotherapy offers the possibility of long-term improvement. Asthma related to dust mites responds well to inhaled preventers, and good environmental control reduces the amount of medication needed.

\section*{Prevention and Self-Care}
\begin{itemize}
    \item Encase mattresses, pillows, and quilts in dust mite-proof covers
    \item Wash bedding weekly in water at 55 degrees Celsius or hotter
    \item Remove carpets from the bedroom and replace with hard flooring where possible
    \item Vacuum regularly using a cleaner with a HEPA filter, and wear a mask while cleaning
    \item Keep soft toys to a minimum, and wash or freeze them regularly
    \item Control humidity with ventilation and dehumidifiers
    \item Keep pets out of the bedroom, as their dander feeds the mites
\end{itemize}

\section*{Sources and Further Reading}
The following reputable sources provide further information for healthcare students and patients seeking reliable, current advice about dust mites and allergy.
\begin{itemize}
    \item National Asthma Council Australia. \textit{Dust mite allergy and asthma control.} https://www.nationalasthma.org.au/
    \item Australasian Society of Clinical Immunology and Allergy. \textit{Dust mite allergy patient information.} https://www.allergy.org.au/
    \item NHS. \textit{Allergies: dust mite allergy.} https://www.nhs.uk/conditions/allergies/
    \item Mayo Clinic. \textit{Dust mite allergy: symptoms and treatment.} https://www.mayoclinic.org/diseases-conditions/dust-mites/
    \item MedlinePlus. \textit{Dust mite allergy.} https://medlineplus.gov/dustmiteallergy.html
    \item American Academy of Allergy, Asthma \& Immunology. \textit{Dust mites patient education.} https://www.aaaai.org/
\end{itemize}
